\documentclass{article}
\usepackage[T1]{fontenc}
\usepackage{lmodern}
\usepackage{graphicx}
\usepackage{amsmath,amssymb}
\usepackage[a4paper,margin=2cm]{geometry}
\usepackage[absolute,overlay]{textpos}
\setlength{\TPHorizModule}{1mm}
\setlength{\TPVertModule}{1mm}
\usepackage[indonesian]{babel}
\usepackage{hyperref}
\setlength{\emergencystretch}{2em}
% unit: Halaman 10

\begin{document}

% === Halaman 10 (llm) ===

\begin{enumerate}
\setcounter{enumi}{3}
\item Bobot total sekarang menjadi $18 - 13 = 5$.
\item Bandingkan 5 dengan bobot terbesar kedua, yaitu 6. Karena $6 > 5$, maka 6 tidak dimasukkan ke dalam \textit{knapsack}. $\rightarrow b_3 = 0$
\item Bandingkan 5 dengan bobot terbesar berikutnya, yaitu 3. Karena $3 \le 5$, maka 3 dimasukkan ke dalam \textit{knapsack}. $\rightarrow b_2 = 1$
\item Bobot total sekarang menjadi $5 - 3 = 2$.
\item Bandingkan 2 dengan bobot terbesar berikutnya, yaitu 2. Karena $2 \le 2$, maka 2 dimasukkan ke dalam \textit{knapsack}. $\rightarrow b_1 = 0$
\item Bobot total sekarang menjadi $2 - 2 = 0$.
\end{enumerate}

\end{document}
